% Complete proof candidate. Individual check status is in result_audit.json.
\section{C.6. Global heights and the eleven rational CM fibres}

We use the equation
\[
 E_t:Y^2=X^3-\frac34t^3(t-1)X+\frac14t^4(t-1)^2,
 \qquad P_t=(0,t^2(t-1)/2).
\]
Its discriminant is $27t^8(t-1)^3$. Our canonical height is
$\widehat h_O(P)=\frac12\lim_{n\to\infty}n^{-2}h(x(nP))$.
In particular the bilinear Mordell--Weil self-pairing used in the
supplied surface calculation is $2\widehat h_O$.
For a finite place $v$ above the rational prime $p$, we normalize
$|p|_v=p^{-1}$ and set $n_v=[K_v:\Q_p]$. At an archimedean
place $n_v$ is one or two, according as the place is real or complex.

\begin{proposition}[ARI C.6.1; Proved: the global formula]\label{ari:C.6.1}
For every number field $K$, every $t\in K\setminus\{0,1\}$,
and the section just specified,
\begin{equation}\label{ari:eq:ch-global}
 \sum_{\sigma:K\hookrightarrow\C}\log N(P_{\sigma(t)})
 =\sum_{v\nmid\infty}n_v\lambda_v(P_t)
          -[K:\Q]\widehat h_O(P_t).
\end{equation}
Here the local heights are invariant under change of Weierstrass
model and satisfy, when $nP\ne O$,
\begin{equation}\label{ari:eq:ch-local-mult}
 \lambda_v(nP)=n^2\lambda_v(P)-\log|\psi_n(P)|_v
                   +\frac{n^2-1}{12}\log|\Delta|_v.
\end{equation}
At a complex embedding $\lambda_v(P)=-\log N(P)$; complex
embeddings are counted separately in \eqref{ari:eq:ch-global}.
For rational $t$, put $A_t=\exp(\sum_p\lambda_p(P_t))$.
Then $A_t$ is a positive real algebraic number and
\[
 N(P_t)=A_t e^{-\widehat h_O(P_t)},\qquad
 D_{\rm part}(t)=\frac6\pi
       \left(\log\frac{A_t}{2\sqrt3}-\widehat h_O(P_t)\right).
\]
Thus $N(P_t)$ is algebraic if and only if
$e^{\widehat h_O(P_t)}$ is algebraic.
\end{proposition}
\begin{proof}
The sigma expression for the archimedean local height is
\[
 \lambda_v(P)=-\log|\sigma(u)|
       +\frac12\Re(u\eta(u))-\frac1{12}\log|\Delta|.
\]
The sigma--theta relation and the Legendre relation identify this
with $-\log N$ exactly, including its constant; this is also the
calculation in Statement C.7.1. The division identity
$\sigma(nu)=(-1)^{n-1}\psi_n(P)\sigma(u)^{n^2}$ gives
\eqref{ari:eq:ch-local-mult}. At finite places the same identity is
the defining multiplication law of the canonical N\'eron function.
Its normalization is uniquely determined by continuity away from
$O$, its logarithmic pole at $O$, and the multiplication law.
These are the three properties printed in
\citet[Section 2.1]{demarco2017}; their duplication law includes
the discriminant term $\frac14\log|\Delta|_v$.
Under a coordinate scaling the changes in the sigma or formal
parameter term and in $\frac1{12}\log|\Delta|_v$ cancel.

Sum the duplication law over places. The product formula cancels
the division polynomial and discriminant terms, so that the sum
$H(P)=[K:\Q]^{-1}\sum_vn_v\lambda_v(P)$ satisfies
$H(2P)=4H(P)$. The pole normalization implies
$H(P)=\frac12h(x(P))+O_E(1)$. Iterating multiplication by two
therefore identifies $H$ with $\widehat h_O$. Moving the
archimedean terms to the other side gives \eqref{ari:eq:ch-global}.
For a finite place, a multiple of $P$ lies in the identity component
of a minimal model. There the local height is a rational multiple
of $\log p$; the multiplication law shows the same for $P$.
Only finitely many places contribute. This proves the assertion
about $A_t$ and both of its consequences.
\end{proof}

\begin{proposition}[ARI C.6.2; Proved: exact finite local-height calculation]\label{ari:C.6.2}
Let $E/\Q$ be a global minimal Weierstrass model. If $mP$ is
in the identity component at $p$, and $mP\ne O$, then
\begin{equation}\label{ari:eq:ch-local-exact}
 \frac{\lambda_p(P)}{\log p}
 =\frac{v_p(\Delta)}{12}
   +\frac{\frac12\max(0,-v_p(x(mP)))-v_p(\psi_m(P))}{m^2}.
\end{equation}
For $m=1$ we read $\psi_1=1$. This formula, with exact rational
coordinates and valuations, gives all the finite contributions
in the tables below.
\end{proposition}
\begin{proof}
On the identity component of a minimal model the canonical
N\'eron function is
$\frac12\log\max(1,|x|_p)-\frac1{12}\log|\Delta|_p$.
This follows either from the formal parameter at $O$ followed
by the group law on the smooth locus, or from the uniqueness
properties used above. Substitute it for $\lambda_p(mP)$ in
\eqref{ari:eq:ch-local-mult} and divide by $m^2\log p$.
The contribution $v_p(\Delta)/(12m^2)$ combines with the
discriminant term to give $v_p(\Delta)/12$.
Membership in the identity component is checked by reduction
to the nonsingular locus of the minimal cubic. Thus the test is
exact also at $2$ and $3$; it does not use a formula restricted
to residue characteristic at least five.
\end{proof}

\begin{proposition}[ARI C.6.3; Proved: the CM models and their height identities]\label{ari:C.6.3}
The eleven rational CM fibres have the global minimal models,
section coordinates, component classes, and finite local heights
in the following tables. On each model $a_1=a_2=0$, and we
write $[a_3,a_4,a_6]$ for
$y^2+a_3y=x^3+a_4x+a_6$. The displayed point is $(0,y_P)$.
The sign of $P$, which is immaterial to all these heights,
is chosen by the recorded change of variables.
The constants $A_t$ in increasing row order are
\[
 \sqrt6,\quad 2^{2/3}3^{1/4},\quad 2^{3/4},\quad
 2,\quad 2,\quad \sqrt3,\quad 1,\quad 3^{5/12},\quad
 1,\quad1,\quad1.
\]
\end{proposition}
\begin{proof}
We substitute the eleven exact values $t=j/1728$, and perform
the integral minimal-model transformations and Tate reduction
algorithm. The transformations and resulting invariants are
printed in \path{code/C_cm_models.out}; the exact valuation
computations are in \path{code/C_cm_local_heights.out}.
For each bad prime the program computes $mP$ by the rational
group law, tests the two partial derivatives of the minimal
cubic, and evaluates $\psi_m$ by its exact recursion. Formula
\eqref{ari:eq:ch-local-exact} then proves each entry. Outside the
listed primes the discriminant is a unit and $P$ is integral,
so the local height is zero. Multiplying $p^{\lambda_p/\log p}$
gives the listed constants $A_t$.

We specify the component labels without relying on an
unrecorded ordering of a resolution graph. The label $0$
means the identity component, and $c_P$ is the class represented
by the explicitly displayed section $P$. At an odd prime of
type $I_0^*$ this nonidentity class is the outer component
indexed by the root zero of the residual Tate cubic; the
scaling $x=pX$ identifies that root. When there are two
rational components, $c_P$ is their unique nonzero class.
The only remaining ambiguity is row 5 at $2$, where the
component group is $(\Z/2)^2$. Here let $T=(-14,0)$;
we label it with the basis $(c_P,c_T)$. Indeed
$P+T=(18,64)$ after the corresponding choice of sign,
and $P,T,P+T$ all reduce to the singular point, whereas
$2P=(25/4,99/8)$ reduces to the identity component.
Thus the three nonzero classes are explicitly separated.

The independent analytic check uses the exact quadratic CM
period ratio, solves $\wp(\mu)=0$ by its convergent Fourier
series, and evaluates the theta product at 270-digit working
precision. The value of $j$ agrees with its prescribed integer
to absolute error less than $10^{-245}$. For every row the
direct value of $\log N$ agrees with
$\log A_t-\widehat h_O(P_t)$ to more than 245 decimal places.
These numerical comparisons test the exact normalization;
the algebraic local-height entries were already proved by
integer and rational arithmetic. The full precision values
are in \path{records/C_cm_theta_pslq.json}.
\end{proof}

\begingroup\footnotesize\setlength{\tabcolsep}{4pt}
\begin{longtable}{r r r r r r}
\toprule Row & $j$ & $a_3$ & $a_4$ & $a_6$ & $y_P$\\\midrule\endhead
1 & 8000 & 0 & -5880 & 153664 & 392\\
2 & 54000 & 0 & -7260 & 234256 & 484\\
3 & 287496 & 0 & -2156 & 38416 & 196\\
4 & 16581375 & 0 & -7732620 & 8275904784 & 90972\\
5 & -3375 & 0 & -140 & 784 & -28\\
6 & -32768 & 1 & -3234 & 72630 & -270\\
7 & -884736 & 1 & -38 & 90 & -10\\
8 & -12288000 & 1 & -1920270 & 1024288020 & -32005\\
9 & -884736000 & 1 & -379260 & 89898842 & -9482\\
10 & -147197952000 & 1 & -347045930 & 2488447882842 & -1577482\\
11 & -262537412640768000 & 1 & -75065352465446420 & 7916029533155605844266842 & -2813543945482\\
\bottomrule\end{longtable}
\begin{longtable}{r r l c c l}
\toprule Row & $p$ & Type & Component & Order & $\lambda_p/\log p$\\\midrule\endhead
1 & 2 & $III^*$ & $c_P$ & 2 & $\frac{1}{2}$\\
1 & 3 & $I_0^*$ & $0$ & 1 & $\frac{1}{2}$\\
1 & 7 & $I_0^*$ & $c_P$ & 2 & $0$\\
2 & 2 & $I_{4}^*$ & $c_P$ & 2 & $\frac{2}{3}$\\
2 & 3 & $III$ & $0$ & 1 & $\frac{1}{4}$\\
2 & 11 & $I_0^*$ & $c_P$ & 2 & $0$\\
3 & 2 & $I_{5}^*$ & $c_P$ & 2 & $\frac{3}{4}$\\
3 & 7 & $I_0^*$ & $c_P$ & 2 & $0$\\
4 & 2 & $I_{8}^*$ & $c_P$ & 2 & $1$\\
4 & 3 & $I_0^*$ & $c_P$ & 2 & $0$\\
4 & 7 & $III$ & $c_P$ & 2 & $0$\\
4 & 19 & $I_0^*$ & $c_P$ & 2 & $0$\\
5 & 2 & $I_{8}^*$ & $c_P$ & 2 & $1$\\
5 & 7 & $III$ & $c_P$ & 2 & $0$\\
6 & 3 & $I_0^*$ & $0$ & 1 & $\frac{1}{2}$\\
6 & 7 & $I_0^*$ & $c_P$ & 2 & $0$\\
6 & 11 & $III$ & $c_P$ & 2 & $0$\\
7 & 19 & $III$ & $c_P$ & 2 & $0$\\
8 & 3 & $IV$ & $0$ & 1 & $\frac{5}{12}$\\
8 & 11 & $I_0^*$ & $c_P$ & 2 & $0$\\
8 & 23 & $I_0^*$ & $c_P$ & 2 & $0$\\
9 & 3 & $I_0^*$ & $c_P$ & 2 & $0$\\
9 & 7 & $I_0^*$ & $c_P$ & 2 & $0$\\
9 & 43 & $III$ & $c_P$ & 2 & $0$\\
10 & 7 & $I_0^*$ & $c_P$ & 2 & $0$\\
10 & 31 & $I_0^*$ & $c_P$ & 2 & $0$\\
10 & 67 & $III$ & $c_P$ & 2 & $0$\\
11 & 7 & $I_0^*$ & $c_P$ & 2 & $0$\\
11 & 11 & $I_0^*$ & $c_P$ & 2 & $0$\\
11 & 19 & $I_0^*$ & $c_P$ & 2 & $0$\\
11 & 127 & $I_0^*$ & $c_P$ & 2 & $0$\\
11 & 163 & $III$ & $c_P$ & 2 & $0$\\
\bottomrule\end{longtable}
\endgroup
\noindent The next table prints 52 significant digits. The last column is both the directly computed $\log N$ and the independently computed $\log A_t-\widehat h_O(P_t)$; the agreement is more precise than the digits printed.
\begingroup\small
\begin{longtable}{r l}
\toprule Row & Height and archimedean logarithm\\\midrule\endhead
1 & $\widehat h_O=0.68444758838339190721754413638128047178526153154571$\\
 & $\log N=0.2114321462306355931886945428090706645762338145457924$\\[3pt]
2 & $\widehat h_O=0.6417998032172362080432024005106271224631326545214062$\\
 & $\log N=0.09495138932308808775043032302545534908240674117445132$\\[3pt]
3 & $\widehat h_O=0.5335302122611885050129127004864161587766141055417204$\\
 & $\log N=-0.01366982684122952294998860939278373271998900477152897$\\[3pt]
4 & $\widehat h_O=0.9940668685687274300038154462317330832788241518412543$\\
 & $\log N=-0.300919688008782120586583324773556515203324017480999$\\[3pt]
5 & $\widehat h_O=0.4673656771814712800075202012628467836356203437033082$\\
 & $\log N=0.2257815033784740294097119201953297844398797906569471$\\[3pt]
6 & $\widehat h_O=0.4286091723813209970260194687110167110834518087280487$\\
 & $\log N=0.120696971952733848671603149750246141240293470183326$\\[3pt]
7 & $\widehat h_O=0.09064583701789963707894077069188164993051948122321821$\\
 & $\log N=-0.09064583701789963707894077069188164993051948122321821$\\[3pt]
8 & $\widehat h_O=0.7366208256267509672496883390906560396314066510934334$\\
 & $\log N=-0.2788657053483719291683361570396036626949522520006211$\\[3pt]
9 & $\widehat h_O=0.6018296198896332298783588845164397135769872973598109$\\
 & $\log N=-0.6018296198896332298783588845164397135769872973598109$\\[3pt]
10 & $\widehat h_O=1.002649286144114870320961646648539589847474352238484$\\
 & $\log N=-1.002649286144114870320961646648539589847474352238484$\\[3pt]
11 & $\widehat h_O=2.165489246796061077105970244492001901750597685962121$\\
 & $\log N=-2.165489246796061077105970244492001901750597685962121$\\[3pt]
\bottomrule\end{longtable}\endgroup


\begin{proposition}[ARI C.6.4; Proved: the maximum as a height inequality]\label{ari:C.6.4}
For $t\in K\setminus\{0,1\}$,
\begin{equation}\label{ari:eq:ch-bound}
 \widehat h_O(P_t)\geq
 \frac1{[K:\Q]}\sum_{v\nmid\infty}n_v\lambda_v(P_t)
                    -\frac12\log2.
\end{equation}
The inequality is strict on this smooth base. Over $\Q$ it
is $\widehat h_O(P_t)\geq\log A_t-\frac12\log2$.
The generic height is $\widehat h_O(P)=1/12$, and the
variation divisor is $D_P=(\infty)/12$. Consequently
\[
 \widehat h_O(P_t)=\frac1{12}h(t)+O(1)
 \qquad(t\in\overline\Q\setminus\{0,1\}),
\]
where the constant is uniform over algebraic $t$ and is for
the fixed family and section.
\end{proposition}
\begin{proof}
Apply the analytic maximum theorem C.8.5 at every embedding
and substitute into \eqref{ari:eq:ch-global}. Equality would
require $j(\sigma(t))=1728$ for every embedding, hence $t=1$,
which has been excluded.

For the generic local heights at $0$ and $1$, the minimal
discriminant orders are $8$ and $3$, and the section meets
the nonidentity component of types $IV^*$ and $III$.
Their local corrections are $2/3$ and $1/4$, respectively,
so their normalized local heights are zero. At infinity,
put $s=1/t$, $X=s^{-2}x$, $Y=s^{-3}y$. The new equation is
\[
 y^2=x^3-\tfrac34(1-s)x+\tfrac14(1-s)^2,
 \qquad P=(0,(1-s)/2),\qquad\Delta=27s(1-s)^3.
\]
At $s=0$ the node is $(1/2,0)$, whereas $P=(0,1/2)$ is
nonsingular. Its height contribution is $1/12$.
At all other function-field places the point is integral
and the discriminant is a unit. This proves
$D_P=(\infty)/12$ and the generic height.

Theorem 1.1 of \citet{demarco2017} applies with $K=\Q$,
$B=\mathbb P^1$, the displayed elliptic surface and rational
section, since the generic height is positive. It equips
this divisor with continuous adelic semipositive metrics
whose associated height is exactly $\widehat h_O(P_t)$.
Comparing with the usual height for $(\infty)/12$ gives the
bounded difference. This is the specialization here of the
variation theorem of Tate, with the local refinements of
Silverman reproduced as Theorems 2.2 and 2.3 in that paper:
outside finitely many places the local correction is zero,
and at each remaining place it extends continuously after
subtracting its prescribed logarithmic singularity.
Our inequality is an explicit archimedean estimate with
constant $\frac12\log2$; the general bounded-variation
statement alone does not supply that constant.

We spell out the comparison with the opened primary author sources.
Theorem 4.1 of \citet[pp.~183--184]{callSilverman1993} concerns a
smooth projective base curve $T$, a degree-one base height $h_T$,
a section, and a fibrewise morphism whose pullback multiplies the
chosen divisor class by a real number $\alpha>1$. It gives
\[
 \lim_{h_T(t)\to\infty}
 \frac{\widehat h_t(P_t)}{h_T(t)}
       =\widehat h_{\Q(T)}(P).
\]
Here take $T=\mathbb P^1$, the class of $(O)$, and the generic
morphism $[2]$, with $[2]^*(O)\sim4(O)$ and $\alpha=4$.
The smooth fibres are precisely the complement of $0,1,\infty$;
the section extends on the regular projective elliptic surface.
All the hypotheses hold, and the limit is $1/12$.
The unnumbered remark immediately following that theorem states
Tate's stronger elliptic-family formula
$\widehat h_t(P_t)=h_{T,D_P}(t)+O_P(1)$, with
$\deg D_P=\widehat h_{\Q(T)}(P)$. This is an author-verified
reproduction of Tate's result, not a claim to have read the
inaccessible original scan \citep{tate1983variation}.
The same distinction applies to \citet{silverman1983specialization}:
the limit theorem is reproduced by the author on slide~22 of
\citet{silverman2006heights}, while slide~21 states the uniform
canonical-versus-Weil-height bound in terms of the base height.
For our explicit bounded remainder and local normalization, the
fully opened Theorems~1.1, 2.2 and 2.3 of
\citet{demarco2017}, applied above, remain the cited proof.
\end{proof}

\begin{proposition}[ARI C.6.5; Proved: rigorous high-precision height enclosures]\label{ari:C.6.5}
For every one of the eleven rational CM parameters in C.6.3,
the files \path{records/C_cm_height_arb.json} and
\path{code/C_cm_height_arb.out} give certified enclosures of
$\log N(P)$ and $\widehat h_O(P)$ of radius less than
$10^{-236}$. These enclosures use no floating-point height
value as a proof step.
\end{proposition}
\begin{proof}
The CM discriminants, in the order of C.6.3, are
\[
 -8,-12,-16,-28,-7,-11,-19,-27,-43,-67,-163.
\]
Enumeration of primitive reduced forms
$(a,b,c)$, $|b|\le a\le c$, $b^2-4ac=D$, with the usual
nonnegative boundary convention, gives exactly one form of
each discriminant, always with $a=1$.
The range is finite and proved by $a\le\sqrt{|D|/3}$.
Thus the CM theorem used in C.5.1 makes the corresponding
singular modulus an integer. The Fourier enclosures below
place it within $1/4$ of the specified $j$, proving its
identity, rather than only numerical closeness. Exact class
polynomials were also computed independently in
\path{code/C_cm_class_polynomials.gp} and its output.

Write $\tau=iT$ for the first four entries and
$\tau=1/2+iT$ for the other seven, with
\[
 (T_1,T_2,T_3,T_4)=(\sqrt2,\sqrt3,2,\sqrt7),\qquad
 T_k=\sqrt{|D_k|}/2\quad(k\ge5).
\]
Put $Q=e^{-2\pi T}$, $q=Q$ or $-Q$ accordingly,
$M=350$, and $Y=2/5$. On $0\le y\le Y$,
\begin{equation}\label{ari:eq:ch-wp-bracket}
 \frac{\wp(1/2-iy)}{\pi^2}
 =\frac1{\cosh^2\pi y}-\frac13+
 8\sum_{n\ge1}\frac{nq^n}{1-q^n}
                 \{1-(-1)^n\cosh(2\pi ny)\}.
\end{equation}
For $a=Qe^{2\pi Y}$, the absolute tail after $M$ in
$\wp$ is at most
\[
 \frac{16\pi^2}{1-Q}
 \frac{a^{M+1}\{M+1-Ma\}}{(1-a)^2}<10^{-875}.
\]
The proof is $1+\cosh(2\pi ny)\le2e^{2\pi nY}$ and
$\sum_{n>M}na^n=a^{M+1}(M+1-Ma)/(1-a)^2$.
Each previously computed decimal supplies only a rational
candidate centre $y_0$. At the two rational endpoints
$y_0\pm10^{-240}$ the truncated expression plus this
infinite tail is respectively strictly positive and strictly
negative. The code records all twenty-two strict signs.
Continuity produces a zero within each bracket. Every zero
of $\wp$ is one of the pair $\pm P$, and $N$ agrees on
that pair, so no unproved choice of root is needed.

The product formula on this bracket has positive real factors:
\begin{align*}
 \log N={}&\log(2\cosh\pi y)-\frac{\pi T}{6}
                      -\frac{\pi y^2}{T}\\
 &+\sum_{n\ge1}\{\log(1+q^ne^{2\pi y})
                         +\log(1+q^ne^{-2\pi y})\}.
\end{align*}
Since $|\log(1+x)|\le |x|/(1-|x|)$ for real $|x|<1$,
the omitted tail is bounded by
\[
 \frac{2e^{2\pi Y}Q^{M+1}}
 {(1-Q)(1-e^{2\pi Y}Q^{M+1})}<10^{-1260}.
\]
All calculations use Arb enclosures at $1152$ bits; the
radius comparisons and endpoint signs are exact interval
comparisons. The resulting radius is less than $10^{-236}$.
Adding the exact finite local logarithms of C.6.3 and
subtracting this enclosure gives the claimed height bound.

For completeness, the singular-modulus certificate uses
$E_4=1+240\sum n^3q^n/(1-q^n)$ and
$\Delta_0=q\prod(1-q^n)^{24}$. Their respective error
bounds, after truncation at $M$, are
\[
 \frac{240(M+1)^3Q^{M+1}(1+4Q+Q^2)}{(1-Q)^5},
 \qquad
 \frac{24Q^{M+1}}{(1-Q)(1-Q^{M+1})}
\]
for $E_4$ and for the logarithm of the product. The first
uses $(M+1+k)^3\le(M+1)^3(k+1)^3$ and the differentiated
geometric series; the second uses the same logarithm bound.
Thus $E_4^3/\Delta_0$ is rigorously enclosed. The complete
certificate is \path{code/C_cm_height_arb.py}, with its
output and rational bracket centres retained.
\end{proof}
